@@M0@@
where @@M1@@ is obtained from @@M2@@ by base change; since this base change is flat, @@M3@@ is a regular immersion, and consequently @@M4@@ is a regular immersion of @@M5@@ into a scheme smooth over @@M6@@, giving the result.

\sgaStableTarget{sga6.e08.prop.1-6}\textbf{Proposition 1.6.} \uline{Let} @@M7@@ \uline{be a morphism and} @@M8@@ \uline{a flat base change}. \uline{If} @@M9@@ \uline{is a complete intersection}, \uline{then the morphism} @@M10@@ \uline{obtained from} @@M11@@ \uline{is one as well}. \uline{The converse is true if} @@M12@@ \uline{is also quasi-compact and surjective}.

One can find a covering of @@M13@@ by open subsets @@M14@@ such that for every @@M15@@, the restriction of @@M16@@ to @@M17@@ factors as @@M18@@, where @@M19@@ is a regular immersion, and @@M20@@ a smooth morphism. The @@M21@@ form an open covering of @@M22@@, the immersions @@M23@@ are regular by VII 1.5, and the morphisms @@M24@@ are smooth, giving the first result. For the converse, EGA IV 2.7.1 (iv) first implies that @@M25@@ is locally of finite presentation, which easily reduces us, using 1.2, to the case where @@M26@@ is a closed immersion, in which case the assertion is precisely VII 1.5.

\sgaStableTarget{sga6.e08.prop.1-7}\textbf{Proposition 1.7.} \uline{Let} @@M27@@ \uline{be a complete-intersection morphism}. \uline{Then} @@M28@@ \uline{is a perfect morphism}.

Since the property is local on @@M29@@, one may work in a neighbourhood of a point @@M30@@ of @@M31@@, and suppose that there exist an open subset @@M32@@ containing @@M33@@, and a factorization:
